% bridg-it.tex

%%%%%%%%%%%%
\begin{frame}{}
	\begin{theorem}
		If $T$, $T'$ are spanning trees of a connected graph $G$
		and \blue{$e \in E(T) \setminus E(T')$},
		then there is an edge \teal{$e' \in E(T') \setminus E(T)$}
		such that \purple{$T - e + e'$} is a spanning tree of $G$.
	\end{theorem}

	\pause
	\vspace{0.30cm}
	\fig{width = 0.50\textwidth}{figs/two-st}
	\begin{center}
		{\small ($T$ is bold, $T'$ is solid; they share two edges.)}
	\end{center}

	\pause
	\vspace{0.30cm}
	\begin{center}
		$T - e$: two components $U$ and $U'$ \\[8pt]
		\pause
		$e'$: an edge across $(U, U')$ in $T'$ \\[8pt]
		\pause
		\purple{$T - e + e'$}: a spanning tree of $G$
	\end{center}
\end{frame}
%%%%%%%%%%%%

%%%%%%%%%%%%%%%%%%%%
\begin{frame}
  \begin{exampleblock}{Bridg-It Game (David Gale, 1958) \qquad
    \href{https://mglaezer.github.io/bridgit/}{\cyan{\bf \large Play with it here!}}}
    \begin{columns}
      \column{0.50\textwidth}
        \fig{width = 0.80\textwidth}{figs/bridg-it}
      \column{0.50\textwidth}
        \uncover<2->{
          \fig{width = 0.90\textwidth}{figs/red-win}
        }
    \end{columns}

    \[
      \red{5 \times 6} \;\text{\it vs.}\; \blue{6 \times 5}
    \]
  \end{exampleblock}
\end{frame}
%%%%%%%%%%%%%%%%%%%%

%%%%%%%%%%%%%%%%%%%%
\begin{frame}
  \fig{width = 0.80\textwidth}{figs/ask-right-questions}
\end{frame}
%%%%%%%%%%%%%%%%%%%%

%%%%%%%%%%%%%%%%%%%%
\begin{frame}
  \begin{center}
    Will Bridg-It \red{\bf end in a tie}?

    \pause
    \vspace{0.80cm}
    No Way!

		\vspace{0.30cm}
		Every path from left to right \blue{crosses}
		every path from top to bottom.
  \end{center}
\end{frame}
%%%%%%%%%%%%%%%%%%%%

%%%%%%%%%%%%%%%%%%%%
\begin{frame}{}
  \begin{center}
    Does \red{\bf Player 2} have a \red{\bf winning strategy}?

    \pause
    \vspace{0.80cm}
    No! By the \blue{\bf strategy stealing argument}.
  \end{center}
\end{frame}
%%%%%%%%%%%%%%%%%%%%

%%%%%%%%%%%%%%%%%%%%
\begin{frame}{}
  \begin{center}
    Does \red{\bf Player 1} have a \red{\bf winning strategy}? \\[8pt]

    \pause
    \vspace{0.80cm}
    Yes!
  \end{center}
\end{frame}
%%%%%%%%%%%%%%%%%%%%

%%%%%%%%%%%%%%%%%%%%
\begin{frame}{}
  \begin{center}
    What are \red{\bf Player 1}'s \red{\bf winning strategies}? \\[8pt]

    \pause
    \vspace{0.80cm}
    Uses \blue{\bf edge-disjoint spanning trees}.
  \end{center}
\end{frame}
%%%%%%%%%%%%%%%%%%%%

%%%%%%%%%%%%%%%%%%%%
\begin{frame}{}
	\begin{center}
    \blue{\small \bf Two edge-disjoint spanning trees
		  to connect all the dots for Player 1}

		\fig{width = 0.60\textwidth}{figs/bridg-it-two-st}

		\pause
		The ``0-th'' turn of Player 2: take away the imaginary (red) edge.

		\fig{width = 0.70\textwidth}{figs/cut-reconnect}

		\pause
		Until \purple{Player 1 has won} or
		when \brown{no single edges remain to be cut}.
	\end{center}
\end{frame}
%%%%%%%%%%%%%%%%%%%%